\section{Introduction}
\label{sec:intro}

Most succinct arguments express a computation as polynomial constraints
over a single finite field and then \emph{arithmetize} every operation
into that field. For bit-oriented computation --- hash functions, block
ciphers, the bitwise glue inside signature schemes --- this is a poor
fit. A $32$-bit $\XOR$, a rotation, a modular addition: each is cheap on
hardware but, forced into a large prime field $\F_q$, expands into many
constraints and many witness elements, because the field has no native
notion of ``the $i$-th bit'' or ``carry into position $j$''. The witness
inflates by one or two orders of magnitude, and since prover cost tracks
witness size, the encoding overhead dominates the proof.

SHA-256 is the canonical stress test. Its compression function is built
almost entirely from three primitives on $32$-bit words --- bitwise
$\XOR$, bitwise $\AND$ (inside $\Ch$ and $\Maj$), and addition modulo
$2^{32}$ --- together with fixed rotations and shifts. A prover that
treats $32$-bit words as opaque field elements must re-derive bit-level
structure with range checks and bit decompositions at every step.

\paragraph{The all-$\Ftwo[X]$ approach.} This note develops, end to end,
an argument that keeps the computation in its native algebra \albert{$\Ftwo[X]$ is not really the native algebra, $\mathbb{Z}[X]$ is}. We
represent each $\wbit$-bit word as a \emph{bit-polynomial} in the cell
ring
\[
  \cellring \;=\; \Ftwo[X]/(X^{\wbit}),
  \qquad \wbit = 32,
\]
a degree-$<\wbit$ polynomial over $\Ftwo$ whose coefficient of $X^b$ is
the word's $b$-th bit. In this representation the bitwise primitives are
ring operations, not emulations:
\begin{itemize}[leftmargin=2em]
  \item \textbf{$\XOR$ is addition} in $\cellring$ ($1+1=0$
    coefficientwise). It costs nothing at the commitment layer: an
    $\XOR$ of committed columns is a committed linear combination, never
    materialised.
  \item \textbf{Shifts and rotations are monomial maps.} $\SHR^c$ and
    $\ROTR^c$ act on a cell as multiplication by $X^{\wbit-c}$ in the
    shift quotient $\Fshift$ or the rotation quotient $\Frot$; they are
    $\Ftwo$-linear, hence realisable as \emph{virtual columns} that are
    not committed but reconstructed from their source.
  \item \textbf{Modular addition is a Binius carry identity.} For a
    binary step $\bp{s} \equiv \bp{a}+\bp{b}\pmod{2^{\wbit}}$, the carry
    word is recovered from a free combination of committed columns, and
    the addition is pinned by a linear identity together with the
    full-adder Hadamard relation in $\Fshift$. Only one carry-out bit per
    row is committed.
  \item \textbf{$\AND$ is a Hadamard product.} Coefficient-wise
    multiplication of two bit-polynomials --- the one genuinely nonlinear
    primitive --- is the only operation that does not commute with the
    linear commitment, and it is the technical heart of the system. \albert{Change wording because we deal with AND and Hadamard not that differently from Binius. Only the merging of Lagrange and monomial basis evaluations may be new (to be checked)}
\end{itemize}
The upshot is an algebraic intermediate representation (AIR) of SHA-256
in which every constraint is linear except for the Hadamard products
that encode $\AND$.

\paragraph{Three layers.} Turning this representation into a SNARK has
three layers, and this paper treats all three.

\emph{(1) Arithmetization} (\Cref{sec:arith}). We lay out the
chained-compression AIR: the column model, the linear identities for
message schedule, $\Sigma$/$\sigma$ mixing, and modular addition, and the
declaration of $\AND$ as a column-level Hadamard relation. The result is
a constraint system over $\Ftwo[X]$ in which rotations and shifts are
virtual and the only ideal-membership constraints capture the
modular-addition carries and word width.

\emph{(2) The binary-field PIOP} (\Cref{sec:piop,sec:hadamard}). We
describe the polynomial interactive oracle reduction that proves the AIR.
Because the trace already lives in $\Ftwo[X]$, the prime-projection step
of the general Zinc\texttt{+} compiler is vacuous; the reduction is an
ideal check, an evaluation projection $\psia$ that substitutes the cell
variable $X=\alpha$ for a random $\alpha\in\K=\mathrm{GF}(2^{128})$, and a
sumcheck over $\K$. The nonlinear Hadamard constraints are read as a
polynomial ideal-membership claim and discharged on the \emph{same}
pipeline --- the same ideal check, the same $\alpha$, the same sumcheck ---
with their operands read off in the Lagrange basis $\psiz$ of the cell, at
the same evaluation point.
Everything reduces to a batch of multilinear-extension evaluation claims at
a common point.

\emph{(3) The commitment} (\Cref{sec:commitment}). The claims are settled
by a hash-based polynomial commitment over $\Ftwo[X]$ whose opening is
\emph{bit-sliced}: it folds the $\{0,1\}$ bit slices of the committed
cells against $\K$-valued evaluation weights to obtain, per column, a
single degree-$<\wbit$ polynomial $a'\in\Kpoly$. The verifier reads the
ordinary evaluation claim off $a'$ via $\psia$ (evaluation at $\alpha$),
and --- crucially for layer (2) --- reads the Hadamard discharge operands
off the \emph{same} $a'$ via $\psiz$ (a subspace-Lagrange projection), so
the nonlinear constraints cost no additional opening.

\paragraph{The unifying thread: $\Ftwo$-linear projections.} A single
idea ties the layers together. Both the evaluation projection $\psia$ and
the Hadamard discharge $\psiz$ are $\Ftwo$-linear maps
$\cellring \to \K$ of the form $w \mapsto \sum_b w_b\,\xi_b$ for a weight
vector $(\xi_b)$ --- monomial powers $\xi_b=\alpha^b$ for $\psia$,
subspace-Lagrange values $\xi_b = L_b(\alpha)$ for $\psiz$ --- both read at the one point $\alpha$. The commitment
binds the full bit-slice fold, so any such functional is a length-$\wbit$
inner product against the one bound object $a'$. The arithmetization is
designed so that its linear constraints survive $\psia$ and its nonlinear
constraints reduce to a $\psiz$ read-off; the PIOP is the reduction that
turns ``the constraints hold'' into ``these projections of $a'$ take
these values''; and the commitment is what makes those projections
checkable against committed data. The same $\Ftwo$-linearity that lets
$\XOR$ be free and shifts be virtual is what lets the random projections
commute through the linear code.

\paragraph{Contributions.} Relative to a prime-field arithmetization of
SHA-256 and to the integer ($\Z[X]$) Zinc\texttt{+} pipeline, the
construction:
\begin{enumerate}[leftmargin=2em]
  \item \emph{commits one bit-slice fold per column, and reads everything
    else off it for free.} $\XOR$, shifts, rotations, and any
    $\Ftwo$-linear combination of them --- including the $\Sigma/\sigma$
    mixing functions --- are virtual: reconstructed from the bound object
    $a'$ rather than committed or argued separately, where \binius{}
    discharges shifts with a dedicated Shift reduction. The same $a'$
    answers \emph{both} projections --- the evaluation $\psia$ for the
    linear constraints and the subspace-Lagrange $\psiz$ for the Hadamard
    constraints --- each a length-$\wbit$ inner product against it, so all
    claims are settled by one opening at one point, with a claim object of
    fixed size $\wbit$ independent of the extension degree;
  \item arithmetizes addition $\bmod\,2^{\wbit}$ with \binius{}'s per-step
    full-adder identity, which the cell ring makes especially cheap: the
    carry word is a free virtual shift, so the prover commits only the
    single carry-out bit per row that the shift cannot supply --- never a
    full $\wbit$-bit carry word;
  \item isolates all nonlinearity into column Hadamard products and
    discharges them, recast as a polynomial ideal-membership claim, on the
    same pipeline as the linear part; and
  \item runs the entire proof over $\K=\mathrm{GF}(2^{128})$ with no prime
    projection, the cell ring never widening through encoding.
\end{enumerate}

\paragraph{Scope and framing.} We present the construction as an
idealized design over generic parameters --- a word width $\wbit$, a
binary extension field $\K$, a linear code --- and record the deployed
instantiation ($\wbit=32$, $\K=\mathrm{GF}(2^{128})$, seven chained
SHA-256 compressions, a rate-$1/4$ $\Ftwo$-RAA code) and its engineering
choices in remarks. The commitment layer is summarised here and
developed in full in the companion note on the bit-sliced $\Ftwo[X]$
opening; the present paper focuses on how arithmetization, PIOP, and
commitment compose into a single argument. We do not treat
zero-knowledge (a future addition) and we are explicit about the
Brakedown-style proof size.

\paragraph{Organisation.}
\Cref{sec:prelim} fixes notation: the cell ring, the field, the AIR
trace model, multilinear extensions, and the two $\Ftwo$-linear
projections.
\Cref{sec:arith} develops the SHA-256 $\Ftwo[X]$ arithmetization.
\Cref{sec:piop} gives the binary-field PIOP (commit, ideal check,
$\psia$, sumcheck) and \Cref{sec:hadamard} the Hadamard discharge.
\Cref{sec:commitment} summarises the bit-sliced commitment and how the
claims are settled.
\Cref{sec:soundness} composes the soundness of the layers under
Fiat--Shamir, and \Cref{sec:impl} reports the deployed instantiation and
its costs.
